\section{Quotient boundaries and bounded height error}
\label{sec:boundary}

We now strengthen Proposition~\ref{prop:uniform-sublinear} to a bounded
error.  A ray quotient records the directions that contract when one
height tends to minus infinity.  The sublinear estimate produces
homeomorphisms of these quotients.  Comparing the volumes of their
images first bounds height variation along $N$-fibers.  Averaging the
image volumes of translates in the source quotient then controls the
remaining $D$-directions.

Boundary actions play a central role in \cite[Theorem~2 and Section~4]{Dymarz2010} and
\cite[Section~4.1]{DymarzFisherXie2025}.  Here the contracting quotient
boundaries are used to compare Haar volumes of boundary images and
bound the height error.

\subsection{Contracting rays and boundary homeomorphisms}

For a generator ray represented by a nonzero form $\alpha$, abbreviate
the objects of Lemma~\ref{lem:cone-quotients} by
\[
 I_\alpha=I_{\R_{\ge0}\alpha},\qquad
 L_\alpha=N/I_\alpha,\qquad p_\alpha(g)=g\cdot1,\qquad
 \ell_\alpha=\ell_{\R_{\ge0}\alpha}.
\]
Multiplying $\alpha$ by a positive scalar does not change these
objects.  Every weight of $\Lie(L_\alpha)$ is a strictly positive
multiple of $\alpha$, and
\begin{equation}
\label{eq:ray-trace}
 \ell_\alpha=c_\alpha\alpha,\qquad c_\alpha>0.
\end{equation}
The positivity follows because at least one generator class on this
ray survives in the quotient abelianization.
Fix one representative of each target generator ray, and use its
pullbacks by the finitely many members of $\mathcal A$ when a source
representative is needed.  We fix left invariant Riemannian metrics
on the finitely many ray quotients.

We use two elementary consequences of the quotient action.
If $y=nd$ is any lift of an element of $G$, then
\begin{equation}
\label{eq:quotient-step}
 p_\beta(yh)=p_\beta(y)\,
     \Ad_d\bigl(p_\beta(h)\bigr).
\end{equation}
Here $\Ad_d$ denotes the induced automorphism on $L_\beta$.
The formula is independent of the lift because $N\cap D\subset I_\beta$.
Also, for each fixed compact set $Q\subset L_\beta$, there are
$C,M,c>0$ such that
\begin{equation}
\label{eq:quotient-contraction}
 \sup_{q\in Q}\|\log(\Ad_d q)\|
 \le C(1+|d|_D)^M e^{c\beta(\pi d)}
 \qquad\text{if }\beta(\pi d)\le0.
\end{equation}
Indeed the quotient has finitely many weights $c_j\beta$, $c_j>0$;
take $c=\min_jc_j$ and use the polynomial bounds on their unipotent
parts from Lemma~\ref{lem:canonical-height},
Equation~\eqref{eq:d-polynomial-unipotent}.  These statements also hold
for compact families of bounded increments.

\begin{lemma}[Boundary maps from sublinear height control]
\label{lem:ray-boundaries}
Let $F$ belong to the fixed Lipschitz quasi-isometry family, let
$A=A_F$, and let $\beta$ be a selected target ray representative.  Put
$\alpha=\beta\circ A$.  There is a homeomorphism
\[
 f_\beta:L_\alpha\longrightarrow L_\beta
\]
and a fixed compact set $Q_\beta\subset L_\beta$ such that
\begin{equation}
\label{eq:boundary-shadow}
 f_\beta(p_\alpha x)\in F(x)\cdot Q_\beta
 \qquad(x\in G).
\end{equation}
The compact sets may be chosen uniformly over the fixed family.
The map obtained from $\bar F$ for the reversed ray pair is
$f_\beta^{-1}$.
\end{lemma}

\begin{proof}
Since $\mathcal A^*$ permutes the generator rays, the source quotient
is defined.  Choose $v\in\mathfrak d$ with $\alpha(\pi v)<0$.
There are only finitely many ray and label choices, so choose all
these $v$'s once.

\smallskip
\noindent\emph{Convergence and a uniform compact shadow.}
For $x_t=x\exp_D(tv)$, Proposition~\ref{prop:uniform-sublinear} gives
\[
 \beta\bigl(\pi F(x_t)-\pi F(x)\bigr)
 =t\alpha(\pi v)+o(t)\le-c_0t
 \qquad(t\ge T_0),
\]
with $c_0,T_0>0$ uniform in $F,x$ and the finite choices.
Set $y_t=F(x)^{-1}F(x_t)$.  Its distance from $1$ is $O(1+t)$.
By Lemma~\ref{lem:path-packing} it has a lift $y_t=n_td_t$ with
$|d_t|_D\le C(1+t)$.  The increments $y_t^{-1}y_{t+1}$ have
uniformly bounded length and hence bounded lifts.  Equations
\eqref{eq:quotient-step} and \eqref{eq:quotient-contraction} imply
\[
 \left\|\log\bigl(p_\beta(y_t)^{-1}p_\beta(y_{t+1})\bigr)\right\|
 \le C(1+t)^M e^{-c_1t}.
\]
The series is summable.  In a left invariant Riemannian metric,
$d(1,\exp Z)\le C\|Z\|$, by using the path $s\mapsto\exp(sZ)$.
Thus the projected trajectory at integer times is Cauchy and stays
in a fixed closed ball.  Completeness and properness of the metric
give convergence in a fixed compact set.
The same estimate for time increments of size at most one gives
convergence for real $t\to\infty$.  The initial interval
$0\le t\le T_0$ contributes a fixed compact set, so the limit lies
in a fixed compact $Q_\beta$.  Equivariance of $p_\beta$ under
the $G$-action therefore gives a limit
\begin{equation}
\label{eq:boundary-limit}
 b(x)=\lim_{t\to\infty}p_\beta F(x\exp_D(tv))
 \in F(x)\cdot Q_\beta .
\end{equation}

\smallskip
\noindent\emph{Independence of the representative.}
We first record the path estimate used for this purpose.
Suppose, for large $t$, a source path stays at distance at most
$C_2t$ from a fixed $x$, has length bounded by a polynomial in $t$,
and satisfies
\[
 \alpha\bigl(\pi(x^{-1}z)\bigr)\le-c_2t
\]
at each of its points $z$.  The uniform sublinear estimate puts
its image, relative to $F(x)$, at $\beta$-heights at most
$-c_2t/2$ for all sufficiently large $t$.
The lifts at all image points have $D$-length $O(t)$, because
their distance from $F(x)$ is $O(t)$.
Subdivide the source path into bounded-length pieces.
Equations~\eqref{eq:quotient-step} and
\eqref{eq:quotient-contraction} bound each projected image increment
by a polynomial times $e^{-ct}$.  There are only polynomially many
pieces, so the projected distance between the endpoints tends to
zero.  Constants here may depend on this fixed collection of paths.

For fixed $d\in D$, connect $x\exp(tv)$ to $xd\exp(tv)$ by
$xd(s)\exp(tv)$, where $d(s)$ is a fixed piecewise smooth path in
$D$ from $1$ to $d$.  Its length is polynomial in $t$, since
conjugation in the nilpotent group $D$ has polynomial norm.
Its points have distance $O(t)$ from $x$, and their relative
$\alpha$-heights equal $t\alpha(\pi v)+O(1)$.
The preceding path estimate gives $b(xd)=b(x)$.

For fixed $u\in I_\alpha$, we will vary the height within the halfspace
$\alpha<0$ to realize $u$ while keeping the $\alpha$-height strongly negative.
Choose a compact convex set $E\subset W$ containing $\pi v$,
entirely in $\alpha<0$, such that every bracket
word containing an outside-ray letter has $h_Z(E)>0$.
Here is why such an $E$ exists.  Every generator form outside the
positive $\alpha$-ray is unbounded above on the halfspace
$\{\alpha\le\alpha(\pi v)\}$.  For each outside form choose a point
in that halfspace on which it is as large as desired, and take
the convex hull together with $\pi v$.  The support maxima for
inside-ray letters remain fixed negative numbers.  Since the
list of words and their lengths is finite, sufficiently large
outside maxima make the support sum of every mixed word positive.

By Lemma~\ref{lem:cone-quotients}, the mixed words span
$\Lie(I_\alpha)$.  Thus the fixed element $u$ belongs to
$B_t(E,O(1))$ for all sufficiently large $t$.
One can see this directly in the all-word description of the box:
the fixed coordinates of $\log u$ in a mixed-word basis are
dominated by its positive exponential widths.
Choose an affine section $\sigma:W\to V$ of $\pi$ with
$\sigma(\pi v)$ equal to the abelianization of $v$.
The compact convex set $E_V=\sigma(E)$ contains that
abelianization and satisfies $\pi E_V=E$.
The paths of Lemma~\ref{lem:support-switch-paths}, applied to $E_V$
and the pivot $\exp(tv)$, connect $\exp(tv)$ to $u\exp(tv)$ in length
$O(t)$ with heights in $t(E+o(1))$.  All their heights remain
strictly negative in $\alpha$ by a linear amount, and all their
points have distance $O(t)$ from $1$.  Left translation by $x$
and the preceding path estimate give $b(xu)=b(x)$.

The stabilizer of $1\in L_\alpha$ is $I_\alpha D$.
The two invariances therefore show that $b(x)$ depends only on
$p_\alpha x$.  Define $f_\beta(p_\alpha x)=b(x)$.
This establishes \eqref{eq:boundary-shadow}.

\smallskip
\noindent\emph{Inverse and continuity.}
By Proposition~\ref{prop:uniform-sublinear},
$A_{\bar F}=A^{-1}$, so the construction for $\bar F$ gives
$\bar f:L_\beta\to L_\alpha$.
For $x_t=x\exp(tv)$ the quotient point $p_\alpha x_t$ is fixed.
Equation~\eqref{eq:boundary-shadow} at $x_t$ supplies representatives
$z_t$ of $f_\beta(p_\alpha x)$ within a fixed distance of $F(x_t)$:
lift the compact set $Q_\beta$ to a bounded set in $G$.
The coarse-inverse estimate and the shadow estimate for $\bar f$
then give
\[
 \bar f(f_\beta(p_\alpha x))\in x_t\cdot Q'
\]
for a fixed compact $Q'\subset L_\alpha$.
The sets on the right contract to $p_\alpha x$, since
$\alpha(\pi v)<0$.  Thus $\bar f f_\beta=\id$.
Interchanging the maps proves $f_\beta\bar f=\id$.

Fix $x$ and $t$.  The projection of a fixed-radius neighborhood
of $x_t$ contains a neighborhood of $p_\alpha x$; this uses only
that the quotient projection is a submersion.
By the Lipschitz bound and \eqref{eq:boundary-shadow}, its image
under $f_\beta$ is contained in $F(x_t)\cdot Q''$, with $Q''$
fixed.  These sets have diameter tending to zero.
Indeed, after normalization by $F(x)$, the $D$-part has length
$O(t)$ and negative linear $\beta$-height, so
\eqref{eq:quotient-contraction} applies; quotient left translations
are isometries.  The fixed action of $F(x)$ is Lipschitz, being
a left translation followed by an automorphism.
This proves continuity at $p_\alpha x$.  The same argument proves
continuity of $\bar f$, completing the proof.
\end{proof}

\subsection{Boundary volumes and variation along \texorpdfstring{$N$}{N}-fibers}

\begin{lemma}[Volumes of boundary images]
\label{lem:boundary-volume}
For each source and target ray pair in Lemma~\ref{lem:ray-boundaries},
there is a fixed compact ball $U\subset L_\alpha$ with nonempty
interior such that
\begin{equation}
\label{eq:boundary-volume}
 \vol_{L_\beta}\bigl(f_\beta(x\cdot U)\bigr)
 \asymp e^{\ell_\beta(\pi F(x))}
 \qquad(x\in G).
\end{equation}
The comparison constants and the choices of balls are uniform
over the fixed quasi-isometry family.
\end{lemma}

\begin{proof}
Fix a unit ball $V_-\subset L_\beta$ and a bounded set of lifts
of radius $R_-$ in $G$.  Denote the common Lipschitz constants
by $L_F,L_{\bar F}$ and the common coarse-inverse error by $E_0$.
For each such lift $b$,
\[
 d_G(\bar F(F(x)b),x)\le L_{\bar F}R_-+E_0.
\]
The shadow estimate for the inverse boundary map therefore puts
$f_\beta^{-1}(F(x)\cdot V_-)$ inside $x\cdot U$, where $U$ is
one compact ball containing
\[
 \{h\cdot q:|h|_G\le L_{\bar F}R_-+E_0,\ q\in Q_{\rm inv}\}.
\]
Here $Q_{\rm inv}$ is its fixed shadow compact.
This choice of $U$ is independent of $x$.

Choose a bounded lift set for this $U$ once, of radius $R_U$.
For every lift $b$ in it,
$d_G(F(xb),F(x))\le L_FR_U$.
The forward shadow estimate gives
$f_\beta(x\cdot U)\subset F(x)\cdot V_+$ for one compact $V_+$.
Thus
\[
 F(x)\cdot V_-\subset f_\beta(x\cdot U)
 \subset F(x)\cdot V_+.
\]
Both comparison sets have positive finite Haar volume.
The $G$-action scales that volume by
$e^{\ell_\beta(\pi F(x))}$, proving the assertion.
There are finitely many ray and label choices, so all constants
can be chosen uniformly.
\end{proof}

\begin{lemma}[Bounded variation on \(N\)-fibers]
\label{lem:leaf-height}
There is a uniform constant $B_N$ such that
\[
 |\pi F(xn)-\pi F(x)|\le B_N
 \qquad(F,\ x\in G,\ n\in N).
\]
In particular $|\pi F(nd)-\pi F(d)|\le B_N$ for all $n\in N,d\in D$.
\end{lemma}

\begin{proof}
Fix a generator $X_i$ and any $t\in\R$.
For every source generator ray other than that of $\alpha_i$,
the element $\exp(tX_i)$ lies in its ideal and acts trivially
on the corresponding quotient.  Hence the two source sets
\[
 x\exp(tX_i)\cdot U,\qquad x\cdot U
\]
are identical.  Lemma~\ref{lem:boundary-volume} bounds the difference
of the associated target trace heights by a uniform constant.
By \eqref{eq:ray-trace}, this bounds every target ray-height
difference except possibly the ray paired with $\alpha_i$.

There is a relation $\sum_\rho c_\rho\rho=0$ over all distinct
target generator rays, with every $c_\rho>0$, by
Lemma~\ref{lem:support-symmetries}.  Applying it to the height
difference bounds the one remaining form as well.
Since the ray forms span $W^*$, there is a uniform constant
$B_1$ with
\[
 |\pi F(x\exp(tX_i))-\pi F(x)|\le B_1
\]
for every $x,i,t$.  No bound on $t$ is imposed.

The sets $B_r(\mathcal C)$ exhaust $N$, since every generator
support maximum on $\mathcal C$ is positive.  By
Lemma~\ref{lem:product-schemes}, every such box element is a
product of at most one fixed number $M$ of generator exponentials.
Their parameters can be arbitrary, because the preceding estimate
is parameter independent.  Thus $B_N=MB_1$ works for every $n$.
Finally $nd=d(d^{-1}nd)$ and $d^{-1}nd\in N$; applying the bound
at the basepoint $d$ proves the last assertion without introducing
any dependence on $d$.
\end{proof}

\subsection{Averaging and bounded height}

Fix a ray pair, its boundary map $f_\beta$, and the ball $U$ from
Lemma~\ref{lem:boundary-volume}.  Pull target Haar measure back to
$L_\alpha$ by this map.  For $d\in D$, the preceding two lemmas make
the mass of every left translate of $d\cdot U$ comparable to
$e^{\ell_\beta(\pi F(d))}$, independently of the translating element.
The next lemma compares this common mass with the source Haar volume
of $d\cdot U$.  Since that volume scales by
$e^{\ell_\alpha(\pi d)}$, the comparison will relate source and target
height characters.

\begin{lemma}[Averaging translates of a Radon measure]
\label{lem:haar-averaging}
Let $L$ be a connected simply connected nilpotent Lie group with
Haar measure $\mu$ and a left invariant Riemannian metric.
Let $\rho$ be a Radon measure on $L$.
Suppose a family of compact positive-Haar-measure sets $E_j$,
including $E_0$, and positive numbers $a_j$ satisfy
\[
 K^{-1}a_j\le\rho(zE_j)\le Ka_j
 \qquad(z\in L,\ j)
\]
for one $K\ge1$.  Then
\[
 \frac{a_j}{\mu(E_j)}\asymp
 \frac{a_0}{\mu(E_0)}
\]
with constants depending only on $K$ and a large-scale volume
doubling constant of $L$, independently of the shapes and diameters
of the sets.  Absolute continuity of $\rho$ is not required.
\end{lemma}

\begin{proof}
For compact $E\subset B(R_E)$ and $T>R_E$, Tonelli's theorem gives
\[
 \int_{B(T)}\rho(zE)\,d\mu(z)
 =\int_L\mu(B(T)\cap yE^{-1})\,d\rho(y).
\]
Nilpotent groups are unimodular, so inversion preserves Haar
measure.  The inner quantity equals $\mu(E)$ for
$y\in B(T-R_E)$, vanishes outside $B(T+R_E)$, and lies between
zero and $\mu(E)$ everywhere.  Consequently
\begin{equation}
\label{eq:haar-kernel-bounds}
 \mu(E)\rho(B(T-R_E))
 \le \int_{B(T)}\rho(zE)\,d\mu(z)
 \le \mu(E)\rho(B(T+R_E)).
\end{equation}

Fix a large-scale doubling constant $D_0$ such that
$\mu(B(2T))\le D_0\mu(B(T))$ for all sufficiently large $T$.
Such a constant follows from the polynomial volume estimates for
nilpotent groups, as in Lemma~\ref{lem:d-scaling}.
Put $q_0=a_0/\mu(E_0)$.
Apply \eqref{eq:haar-kernel-bounds} to $E_0$, first with averaging
radius $T-R_{E_0}$ and then with radius $T+R_{E_0}$.
The assumed translate bounds and doubling give, for sufficiently
large $T$,
\[
 (KD_0)^{-1}q_0\mu(B(T))
 \le\rho(B(T))\le KD_0q_0\mu(B(T)).
\]
For a fixed $j$, choose $T$ so large that these bounds apply to
$T-R_{E_j}$ and $T+R_{E_j}$ and that $T\ge2R_{E_j}$.
Use \eqref{eq:haar-kernel-bounds} for $E_j$ and the assumed
translate bounds once more.  Doubling yields
\[
 (K^2D_0^2)^{-1}q_0
 \le\frac{a_j}{\mu(E_j)}
 \le K^2D_0^2q_0.
\]
The required averaging radius may depend on $j$; the displayed
constants do not.  Only integration of nonnegative measurable
functions was used, so the proof applies equally to singular
Radon measures.
\end{proof}

\begin{proof}[Proof of Theorem~\ref{thm:bounded-height}]
If $W=0$, the assertion is immediate.  Otherwise first take a
Lipschitz pair from the uniform family of
Lemma~\ref{lem:lipschitz-pairs}, with label $A_F$ from
Proposition~\ref{prop:uniform-sublinear}.
Fix a target ray $\beta$ and its source ray
$\alpha=\beta\circ A_F$.  Let $f_\beta$ and $U$ be given by
Lemmas~\ref{lem:ray-boundaries} and \ref{lem:boundary-volume}.
Define a Radon measure on $L_\alpha$ by
\[
 \rho(E)=\vol_{L_\beta}(f_\beta(E)).
\]
This is a measure because $f_\beta$ is a homeomorphism; no
absolute-continuity property is asserted.

For $d\in D$, put $U_d=d\cdot U$ and
$a_d=e^{\ell_\beta(\pi F(d))}$.
Every $z\in L_\alpha$ has a representative $n\in N$, and
$zU_d=nd\cdot U$.
Lemmas~\ref{lem:boundary-volume} and \ref{lem:leaf-height} give
\[
 \rho(zU_d)\asymp a_d
 \qquad(z\in L_\alpha,d\in D)
\]
with one uniform comparison constant.
Lemma~\ref{lem:haar-averaging}, applied to these sets and the
reference set $U_1=U$, implies
\[
 \frac{e^{\ell_\beta(\pi F(d))}}{\vol_{L_\alpha}(U_d)}
 \asymp
 \frac{e^{\ell_\beta(\pi F(1))}}{\vol_{L_\alpha}(U)}.
\]
The source action scales Haar by $e^{\ell_\alpha(\pi d)}$.
Taking logarithms therefore gives
\begin{equation}
\label{eq:trace-height-bound}
 \ell_\beta(\pi F(d)-\pi F(1))
 =\ell_\alpha(\pi d)+O(1)
 \qquad(d\in D),
\end{equation}
with a uniform error.

We identify its linear term using the already proved sublinear
estimate.  For $w\in W$, choose $H\in\mathfrak d$ with $\pi H=w$
and set $d_t=\exp_D(tH)$.  Its length is $O(t)$, so
\eqref{eq:trace-height-bound} and
Proposition~\ref{prop:uniform-sublinear} give
\[
 t\ell_\alpha(w)+O(1)
 =\ell_\beta(\pi F(d_t)-\pi F(1))
 =t\ell_\beta(A_Fw)+o(t).
\]
After division by $t$,
$\ell_\alpha=\ell_\beta\circ A_F$.
The target ray traces span $W^*$ by \eqref{eq:ray-trace}.
A fixed spanning subfamily of \eqref{eq:trace-height-bound}
therefore yields
\[
 |\pi F(d)-\pi F(1)-A_F\pi d|\le B_D
 \qquad(d\in D),
\]
with uniform $B_D$.
For any $g=nd$, Lemma~\ref{lem:leaf-height} and $\pi n=0$
extend this bound to $g$, with bound $B_D+B_N$.
All finite ray and label choices were absorbed into the constants,
so the conclusion holds for the whole fixed family.

Finally an original $(K,C)$ quasi-isometry is at a uniformly
bounded distance from its Lipschitz approximation by
Lemma~\ref{lem:lipschitz-pairs}.  Since $\pi$ is Lipschitz,
changing back alters the error by a uniform constant, proving
the stated theorem for the original maps.
\end{proof}
